% =====================================================================
%  Scale-Aware Joint Compression -- Methodology
%  Drop-in section for a NeurIPS/ICML/ACL-style manuscript.
%
%  Required preamble packages:
%    \usepackage{booktabs, amsmath, amssymb, algorithm, algpseudocode}
%    \usepackage{multirow, siunitx}
%
%  Every quantity below is traceable to the repository:
%    method   : docs/method_definition.md, src/.../compression/{arms,layerwise,reconstruct}.py
%    protocol : docs/protocol_freeze.md, docs/protocol_amendment_a1.md
%    evidence : docs/findings_log.md (F-05, F-37, F-38, F-39--F-41, F-44)
% =====================================================================

\section{Methodology}
\label{sec:method}

\subsection{Study design and scope}
\label{subsec:scope}

We measure a single design choice under controlled conditions: whether interleaving
pruning and quantisation within one optimisation objective (\emph{joint}) outperforms
applying them in sequence (\emph{sequential}), and whether any such advantage varies with
model scale. We introduce no new compression primitive. Both the pruning criterion and the
quantiser are deliberately standard, so that the measured quantity isolates the
\emph{pipeline ordering} rather than the strength of an underlying heuristic; a bespoke
criterion would confound ``joint helps'' with ``our criterion helps''.

The comparison is therefore between one specific sequential implementation and one specific
joint implementation, and results should be read as evidence about these pipelines rather
than about joint compression as a family. Concretely, we evaluate five arms:

\begin{itemize}
  \item \textbf{Dense} (FP32): the reference against which all retention is measured.
  \item \textbf{Pruning-only}: mask, then reconstruct. Weights remain FP32, which makes this
        the only arm eligible for native CPU latency measurement (\S\ref{subsec:deployment}).
  \item \textbf{Quantisation-only}: no mask; reconstruct under the quantisation grid.
  \item \textbf{Sequential} ($P\!\rightarrow\!Q$ and $Q\!\rightarrow\!P$): the two orderings.
  \item \textbf{Joint}: alternating layerwise co-optimisation of mask and quantiser.
\end{itemize}

\paragraph{Method family.} The unit of optimisation is \emph{local reconstruction steps per
layer}, not global optimiser steps: we perform layerwise post-training reconstruction rather
than quantisation-aware fine-tuning. This choice is a memory-compute constraint made
explicit --- full-model fine-tuning of a $10^9$-parameter model does not fit within the
\SI{6}{\giga\byte} device budget available to this study, whereas layerwise reconstruction
requires only a single $(d_{\mathrm{in}}\!\times\!d_{\mathrm{in}})$ Gram buffer per layer
(\S\ref{subsec:objective}). A short full-model recovery phase is examined separately as a
post-hoc ablation (\S\ref{subsec:recovery}) and is not part of the primary comparison.

\subsection{The layerwise reconstruction objective}
\label{subsec:objective}

Let $W \in \mathbb{R}^{d_{\mathrm{out}} \times d_{\mathrm{in}}}$ denote the weight of a
targeted linear layer and $X \in \mathbb{R}^{n \times d_{\mathrm{in}}}$ the calibration
activations observed at its input. Let $M \in \{0,1\}^{d_{\mathrm{out}} \times d_{\mathrm{in}}}$
be a binary keep-mask and $Q_b(\cdot)$ a symmetric uniform quantiser at bit width $b$. Since
\texttt{nn.Linear} computes $Y = XW^{\top}$, the layerwise reconstruction objective is

\begin{align}
  \mathcal{L}_{\mathrm{rec}}(M, Q_b)
    &= \bigl\| XW^{\top} - X\,(M \odot Q_b(W))^{\top} \bigr\|_F^{2}
    \label{eq:objective} \\[2pt]
  \text{subject to}\quad
    &\operatorname{sparsity}(M) = s, \qquad \text{bit width}(Q_b) = b .
    \nonumber
\end{align}

Expanding \eqref{eq:objective} over output channels yields a sum of quadratic forms in the
Gram matrix $H = X^{\top}X$:

\begin{align}
  \mathcal{L}_{\mathrm{rec}}
    = \sum_{o=1}^{d_{\mathrm{out}}}
      (w_o - \hat{w}_o)^{\top} H \, (w_o - \hat{w}_o),
  \qquad
  \hat{w}_o = m_o \odot Q_b(w_o),
  \label{eq:gram}
\end{align}

where $w_o$ and $m_o$ denote the $o$-th rows of $W$ and $M$. Two properties of
\eqref{eq:gram} are load-bearing for this study. First, the objective depends on the
calibration data \emph{only} through $H$, so memory cost is fixed at one
$(d_{\mathrm{in}}\!\times\!d_{\mathrm{in}})$ buffer per layer irrespective of calibration
set size. Second, the right-hand side of the induced normal equations is formed over the
\emph{full} dense row, including entries about to be zeroed; surviving weights therefore
absorb the contribution of pruned weights rather than the layer simply losing it. This error
compensation accounts for most of the benefit of reconstruction over naive masking.

Layers are processed one at a time in depth order, with calibration activations propagated
through the \emph{already-compressed} prefix, so that accumulated error is realistic rather
than optimistic.

\paragraph{Solver.} We use a damped alternating refinement solver. Given a fixed mask, the
solver alternates between a continuous correction against the current residual and a
projection back onto the quantisation grid, retaining a candidate only when the objective
strictly improves:

\begin{align}
  \hat{W}^{(t+1)} =
  \begin{cases}
    \Pi_{M, Q_b}\!\bigl(\hat{W}^{(t)} + \Delta^{(t)}\bigr)
      & \text{if } \mathcal{L}_{\mathrm{rec}}\bigl(\Pi_{M, Q_b}(\hat{W}^{(t)} + \Delta^{(t)})\bigr)
        < \mathcal{L}_{\mathrm{rec}}\bigl(\hat{W}^{(t)}\bigr), \\[4pt]
    \hat{W}^{(t)} & \text{otherwise,}
  \end{cases}
  \label{eq:accept}
\end{align}

where $\Pi_{M,Q_b}(\cdot) = M \odot Q_b(\cdot)$ and $\Delta^{(t)}$ solves the ridge-damped
system $(H + \lambda \bar{h} I)\Delta = H(W - \hat{W}^{(t)})$ with $\bar{h}$ the mean Gram
diagonal and $\lambda = 10^{-2}$. Iterate zero is the naive baseline (mask and
round-to-nearest). The accept-only-if-better rule in \eqref{eq:accept} is necessary rather
than cosmetic: projection onto a discrete grid is not guaranteed to reduce a quadratic
objective, so an unguarded loop can terminate worse than it began.

\subsection{Targeted modules}
\label{subsec:modules}

All arms compress an identical module set, resolved through a single code path. An arm that
compressed more layers than another would exhibit a ``gain'' that is in fact a coverage
difference. We target linear layers inside decoder blocks only:

\begin{itemize}\setlength{\itemsep}{2pt}
  \item \textbf{GPT-NeoX / Pythia}: \texttt{attention.query\_key\_value} (fused QKV) and
        \texttt{attention.dense}; \texttt{mlp.dense\_h\_to\_4h} and
        \texttt{mlp.dense\_4h\_to\_h}.
  \item \textbf{Qwen2}: \texttt{self\_attn.\{q,k,v,o\}\_proj};
        \texttt{mlp.\{gate,up,down\}\_proj}.
\end{itemize}

Input embeddings, positional embeddings, the output head, LayerNorm parameters and all
biases are excluded. The exclusion of embeddings is required for the scale question to be
well posed: embedding parameters constitute a large fraction of a 160M model and a much
smaller fraction of a 1B model, so including them would make the \emph{effective}
compression budget vary with scale, rendering a scale trend in joint gain inseparable from a
scale trend in how much of the model was compressed. Consequently, reported compression
ratios are ratios over compressible parameters, and the scale axis is \textbf{targeted
non-embedding parameter count}. For Qwen2.5 the exclusion is additionally forced by tied
embeddings: compressing \texttt{lm\_head} would compress the input embedding as a side
effect.

\subsection{Pruning and quantisation primitives}
\label{subsec:primitives}

\paragraph{Saliency.} We score weights by activation-weighted magnitude,

\begin{align}
  S_{ij} = \lvert W_{ij} \rvert \cdot \lVert X_{:,j} \rVert_2 ,
  \label{eq:saliency}
\end{align}

so that a weight is deemed unimportant only if it is both small \emph{and} multiplies a
low-energy input. The column norms in \eqref{eq:saliency} are already accumulated for the
reconstruction solver, so the criterion is obtained at no additional cost. Pruning is
unstructured at a fixed global sparsity $s$, held constant across model sizes. Masks are
realised exactly and verified on the converted, reloaded artefact rather than in memory.

\paragraph{Quantisation.} We use weight-only symmetric quantisation, per-channel by default
with per-group granularity (group size $128$) for sub-8-bit schemes, at $b \in \{8, 4\}$.
Activations remain in floating point. The scheme, granularity, group size and excluded
modules are identical across all arms by construction: every arm includes the same
configuration fragment rather than restating it.

We distinguish \emph{fake} quantisation (FP32 storage, values snapped to the grid), which
the solver operates on, from \emph{real} quantisation (packed low-bit storage plus scales).
Only the converted artefact is measured for size or latency; a model that was fake-quantised
but never converted yields correct-looking quality alongside a meaningless size. Run records
carry \texttt{is\_converted} and \texttt{storage\_efficiency} so that this failure mode is
detectable after the fact.

\subsection{Arm definitions}
\label{subsec:arms}

The five arms are thin declarations over one shared driver; the only quantity an arm may
vary is the order in which it calls the solver. This is what makes ``what qualifies as
joint'' checkable in code rather than by inspection. Writing $\textsc{Solve}(M, b)$ for one
call of the solver of \S\ref{subsec:objective} against the \emph{dense} weight under mask
$M$ at bit width $b$, and $\textsc{Mask}(\cdot)$ for the top-$(1{-}s)$ selection under
\eqref{eq:saliency}:

\begin{align}
  \textbf{Pruning:}\quad
    & M \leftarrow \textsc{Mask}(W), \quad
      \hat{W} \leftarrow \textsc{Solve}(M, \varnothing)
    \nonumber \\[2pt]
  \textbf{Quantisation:}\quad
    & M \leftarrow \mathbf{1}, \quad
      \hat{W} \leftarrow \textsc{Solve}(M, b)
    \nonumber \\[2pt]
  \textbf{Sequential } P\!\rightarrow\!Q:\quad
    & M \leftarrow \textsc{Mask}(W), \quad
      \tilde{W} \leftarrow \textsc{Solve}(M, \varnothing),
    \nonumber \\
    & \hat{W} \leftarrow \textsc{Solve}(M, b \mid \Sigma \text{ fitted on } \tilde{W})
    \nonumber \\[2pt]
  \textbf{Sequential } Q\!\rightarrow\!P:\quad
    & \tilde{W} \leftarrow \textsc{Solve}(\mathbf{1}, b), \quad
      M \leftarrow \textsc{Mask}_{Q}(\tilde{W}),
    \nonumber \\
    & \hat{W} \leftarrow \textsc{Solve}(M, b \mid \Sigma \text{ reused exactly})
    \nonumber
\end{align}

A critical implementation detail: \emph{every} arm reconstructs against the dense weight, so
the objective in \eqref{eq:objective} is identical across arms. An earlier revision passed
each arm's intermediate weight as the reconstruction target, causing $P\!\rightarrow\!Q$ to
minimise distance to its pruned approximation and $Q\!\rightarrow\!P$ to minimise distance
to its quantised one, while the joint arm minimised distance to dense. The arms were then
optimising three different objectives, and the asymmetry inflated joint gain --- precisely
the failure that matched-condition fairness exists to prevent. What an arm \emph{may} vary
is the provenance of the quantisation grid, which is a genuine pipeline property:
$P\!\rightarrow\!Q$ fits the grid on the post-pruning distribution; $Q\!\rightarrow\!P$
fits once on the dense tensor and reuses it \emph{exactly}, since refitting after the mask
moves would convert it into a second joint arm.

\subsection{The joint arm}
\label{subsec:joint}

The joint arm performs alternating layerwise co-optimisation of the mask and the quantiser
(Algorithm~\ref{alg:joint}). Two conditions define membership in the joint family, and both
are asserted by regression tests that fail if the arm degenerates into ``prune fully, then
apply ordinary PTQ'': (i) mask decisions are evaluated under fake-quantised weights, and
(ii) quantisation scales are re-estimated after the mask changes.

\begin{algorithm}[t]
\caption{Alternating layerwise co-optimisation (joint arm), per layer}
\label{alg:joint}
\begin{algorithmic}[1]
  \Require dense weight $W$, Gram $H = X^\top X$, sparsity $s$, bit width $b$, iterations $K$
  \State $M^{(0)} \gets \textsc{Mask}(W)$
      \Comment{first mask has no quantiser to see}
  \State $\mathcal{I} \gets \textsc{Solve}(M^{(0)}, b)$
      \Comment{incumbent: weight, codes, scales, loss}
  \State $M^{\star} \gets M^{(0)}$
  \For{$k = 1$ \textbf{to} $K-1$}
    \State $\Sigma \gets \textsc{FitScales}(\mathcal{I}.w, M^{\star})$
        \Comment{re-estimation requirement}
    \State $M^{(k)} \gets \textsc{Mask}_{Q}\!\bigl(\textsc{FakeQuant}(\mathcal{I}.w, \Sigma)\bigr)$
        \Comment{mask scored under quantised weights}
    \State $\mathcal{P} \gets \textsc{Solve}\bigl(M^{(k)}, b \mid \text{grid fitted on } \mathcal{I}.w\bigr)$
    \If{$\mathcal{P}.\ell < \mathcal{I}.\ell$}
        \Comment{incumbent guard}
      \State $\mathcal{I} \gets \mathcal{P}$; \quad $M^{\star} \gets M^{(k)}$
    \EndIf
    \State record $\bigl(k,\; \mathcal{I}.\ell,\; \mathcal{P}.\ell,\; \text{accepted},\;
           \lVert M^{(k)} \neq M^{\star} \rVert_1 / \lvert M \rvert\bigr)$
  \EndFor
  \State \Return $(M^{\star}, \mathcal{I})$
\end{algorithmic}
\end{algorithm}

\paragraph{The incumbent guard.} The solver's own accept-only-if-better rule
\eqref{eq:accept} protects a \emph{fixed} mask; nothing protects the outer loop across mask
\emph{changes}. Without the guard on line~8, an iteration selecting a worse mask replaces a
better feasible solution the search had already found. This is empirically measurable: at
$30\%$ sparsity with $b{=}4$ on Pythia-160M, four alternations scored worse than two
($73.17$ against $71.87$ perplexity), the signature of a wandering rather than converging
search. The guard does not tilt the comparison toward joint --- it prevents an alternating
procedure from discarding a solution it had already reached, and leaves the step count
unchanged.

\paragraph{Mask scoring under quantisation.} The joint mask is scored on quantised weights,
$S_{ij} = \lvert Q_b(W_{ij}) \rvert \cdot \lVert X_{:,j} \rVert_2$, optionally replaced by a
keep-versus-prune benefit criterion. A mask ranked on untouched FP32 weights is chosen in
ignorance of the grid, which is exactly the degenerate case the definition excludes. The
consequence of this design is that the joint mechanism is \emph{precision-dependent by
construction}, a prediction we state before presenting results: at $b{=}8$ quantisation error
is small enough that the saliency ranking survives largely intact, so joint and sequential
masks nearly coincide and the arms should be near-indistinguishable. The $b{=}8$ budget is
therefore a \emph{control}, not a second treatment.

\subsection{Fairness constraints}
\label{subsec:fairness}

The dominant threat to validity in a comparison of this form is that the joint arm receives
more optimisation than the sequential arm and the resulting quality difference is attributed
to the method. We enforce matching in code, and assert it both before and after each run.

\begin{table}[t]
\centering
\caption{Matched conditions and their enforcement mechanism. Assertions run from the
configuration before a sweep begins (hours of compute are at stake) and again from the
persisted records afterwards.}
\label{tab:fairness}
\begin{tabular}{@{}lll@{}}
\toprule
\textbf{Quantity} & \textbf{Requirement} & \textbf{Enforcement} \\
\midrule
Target sparsity        & identical            & shared config fragment; asserted \\
Bit width, granularity & identical            & shared config fragment; asserted \\
Targeted modules       & identical            & single \texttt{select\_compressible\_modules} path \\
Layer processing order & identical            & shared layerwise driver, depth order \\
Calibration tensors    & byte-identical       & indices from fixed \texttt{calibration\_seed} \\
Solver passes / layer  & equal across arms    & \texttt{assert\_arms\_can\_be\_matched} (pre-run) \\
Total local steps      & equal across model   & \texttt{assert\_matched\_plans} (post-run) \\
Artefact format        & identical            & single conversion path; \texttt{storage\_efficiency} \\
Evaluation window      & identical            & dataset fingerprint compared per record \\
\bottomrule
\end{tabular}
\end{table}

Two mechanisms deserve emphasis. First, the sequential $P\!\rightarrow\!Q$ pipeline performs
a \emph{second} reconstruction pass specifically so that its solver-call total equals the
joint arm's $K$ alternations; without it, the joint arm would receive strictly more
optimisation. With the sweep solver each call is one deterministic pass, so \emph{passes},
not nominal step counts, are the fairness unit, and we set $K = 2$ to match. We note that an
earlier screening grid ran with $K = 4$ against sequential's two passes --- i.e.\ the joint
arm on twice the optimisation budget --- which motivated promoting this check from a
documented convention to an executable pre-run assertion.

Second, calibration indices derive from a dedicated \texttt{calibration\_seed} rather than
the run seed, so that varying the replicate index changes the calibration draw (and hence
$H$, the mask, and the solution) without perturbing anything else, and both arms of a pair
share that draw exactly.

\paragraph{Acknowledged limitation.} Equal steps is not equal difficulty. The joint arm
solves a harder problem within the same budget, so a matched comparison may understate what
joint could achieve under a budget tuned for it, just as an unmatched comparison overstates
it. The matched comparison is the defensible one, and we report this asymmetry rather than
resolving it.

\subsection{Block-wise GPU offloading and dependency-group capture}
\label{subsec:offload}

Two implementation properties of the driver materially affect what can be run and what the
numbers mean.

\paragraph{Block-wise offloading.} Holding a full $10^{9}$-parameter model resident on a
\SI{6}{\giga\byte} device causes the compression stage to peak at \SI{6.31}{\gibi\byte},
completing only by spilling to host memory at roughly $7\times$ the solve time. We instead
hold \emph{one decoder block} on the device at a time, keeping the remainder on the host, at
a cost of two host--device transfers per block. This changes no reported number --- the
resident and offloaded paths were verified bit-identical --- but it is what makes the 1B
scale point runnable at all, reducing peak reservation to \SI{4.29}{\gibi\byte} with no
spill.

\paragraph{Capture per dependency group.} Activations are captured once per
\emph{dependency group} rather than once per block. A module whose input is produced by
another module in the same block must be fitted against the inputs it will actually receive
--- that is, after its producer has been compressed. Capturing once per block would fit an
MLP down-projection against activations produced by the \emph{dense} up-projection, inputs
that never occur at inference; that is block-wise reconstruction described as layerwise. In
Pythia, attention output and MLP projections form separate groups; Qwen2's sequential
residual structure makes the MLP depend on the whole block, yielding four groups.

\subsection{Experimental grid}
\label{subsec:grid}

\begin{table}[t]
\centering
\caption{Model suite. The scale axis is targeted (non-embedding) parameter count. Depth is
\emph{not} monotone in the Pythia suite: pythia-1b is shallower than pythia-410m while being
$2.7\times$ wider in targeted parameters. Qwen2.5-0.5B is an external-validity check, not a
scale point.}
\label{tab:models}
\begin{tabular}{@{}llrrrl@{}}
\toprule
\textbf{Model} & \textbf{Revision} & \textbf{Blocks} & \textbf{Targeted params} & \textbf{Modules} & \textbf{Role} \\
\midrule
pythia-160m    & \texttt{50f5173d} & 12 & \num{84934656}  & 48 & scale sweep \\
pythia-410m    & \texttt{dd47b0e}  & 24 & \num{301989888} & 96 & scale sweep \\
pythia-1b      & \texttt{f73d7dcc} & 16 & \num{805306368} & 64 & scale sweep \\
\midrule
qwen2.5-0.5b   & \texttt{060db649} & 24 & ---             & 96 & external validation \\
\bottomrule
\end{tabular}
\end{table}

\paragraph{Budgets.} Two budgets are run per model, varying \emph{precision} at fixed
sparsity so that the precision-dependence prediction of \S\ref{subsec:joint} is directly
testable:

\begin{itemize}
  \item \textbf{Moderate} --- $30\%$ sparsity, $b{=}8$. The control condition.
  \item \textbf{Aggressive} --- $30\%$ sparsity, $b{=}4$. The headline condition.
\end{itemize}

Both budgets were selected on the validation split during a screening phase, under a
pre-declared eligibility rule (technically stable, measurably but not catastrophically
degraded), and frozen before the test split was accessed. The sequential order was likewise
frozen per (model, budget) cell on validation, so only one order is run on test.

\paragraph{Data and evaluation.} Calibration draws $128$ sequences of $512$ tokens from the
WikiText-2 training split; evaluation uses $512$ non-overlapping windows of $512$ tokens.
Calibration is disjoint from evaluation by construction, and overlap raises rather than
warns. The quality metric is perplexity retention relative to each model's own dense
baseline,

\begin{align}
  R = 100 \cdot \frac{\mathrm{PPL}_{\mathrm{dense}}}{\mathrm{PPL}_{\mathrm{compressed}}},
  \label{eq:retention}
\end{align}

rather than raw perplexity, which falls with scale for reasons unrelated to compression and
would render any cross-scale trend uninterpretable. The dependent variable is \emph{joint
gain},

\begin{align}
  G = R_{\mathrm{joint}} - R_{\mathrm{sequential}},
  \label{eq:gain}
\end{align}

in percentage points, at a matched compression budget, so that $G > 0$ always denotes a
joint advantage.

\paragraph{Split discipline.} The validation split is a declared \emph{selection surface}:
budgets, sequential orders, and every exploratory estimate live there. The test split was
accessed once, under a frozen configuration, with no tuning afterwards. Every record stores
its evaluation split and dataset fingerprint, and pairing keys include both --- a safeguard
added after a defect was found in which a validation record could be paired against a
test record because experiment identifiers did not encode the split.

\subsection{Statistical analysis}
\label{subsec:stats}

Replication is over \emph{paired calibration draws}: replicate $r$ selects a fixed
calibration subset shared by both arms of a cell, so the joint and sequential arms of a pair
differ only in pipeline. The compression procedure itself is deterministic given a draw,
so treating replicates as ordinary run-seed repetitions would produce duplicates. We use
$R = 8$ paired replicates at 160M and 410M and $R = 5$ at 1B, the latter a schedule
concession whose consequence we state explicitly: at $R = 5$ the smallest attainable
two-sided exact sign-test $p$-value is $0.0625$, so the 1B leg contributes effect sizes and
sign consistency but cannot reach conventional significance at any effect size. This is a
design limit, not a finding.

Three analyses are pre-specified:

\begin{enumerate}
  \item \textbf{Exact two-sided sign test} over the $R$ paired gains, testing sign
        consistency without distributional assumptions.
  \item \textbf{Paired block bootstrap} over whole evaluation windows ($10^4$ resamples, one
        index draw applied to both arms), reported on per-token NLL advantage. Windows rather
        than tokens are the resampling unit because neighbouring tokens are dependent and a
        token-level bootstrap understates the interval.
  \item \textbf{Holm--Bonferroni correction} across the six (scale $\times$ budget) cells
        examined. Reporting the largest of six raw $p$-values without correction is precisely
        the inflation our pre-registered threat analysis anticipated; corrected values are the
        ones we quote.
\end{enumerate}

\paragraph{Pre-registered practical-importance criterion.} A cell is declared practically
important only if it satisfies \emph{both} conditions: mean gain $\bar{G} \geq
\SI{1.0}{pp}$ \emph{and} $\operatorname{sign}(G_r)$ consistent across all $r \in
\{1,\dots,R\}$. The threshold was fixed before any test-split data existed and is not
adjusted post hoc. We report near-misses against it explicitly rather than renegotiating the
bar.

\paragraph{Provenance gate.} Before any test-split number is read, an automated audit
verifies that every planned cell exists, is complete, non-stale, produced under a single
\texttt{method\_version}, and originates from a single host. The confirmatory grid comprises
$171$ records forming $42$ joint/sequential pairs; the audit must report zero failures,
zero stale and zero missing records before analysis proceeds.

\subsection{Deployment measurement protocol}
\label{subsec:deployment}

Deployment quantities --- latency, throughput, peak memory and checkpoint size --- are
measured exclusively on CPU, on a single designated benchmark host, with a pinned thread
count and fixed batch and sequence lengths. Compression and quality evaluation are portable
across accelerators (they differ only by floating-point reduction order, on the order of
$10^{-5}$ relative, against effects of order $10^{-2}$); a latency, by contrast, is a
property of a machine. Accordingly no results table combines records from more than one
host, and every record carries full hardware metadata so that a violation is detectable
after the fact.

We report prefill and decode latency separately (median and IQR over $60$ samples per cell,
with model-order rotation so thermal drift is spread across arms rather than loaded onto
whichever ran first), because a blended generation latency conceals the compute-bound versus
bandwidth-bound distinction the sparsity question turns on.

\paragraph{Scope restriction on $b{=}4$.} The CPU quantisation backend provides a mature
INT8 path but no equally mature 4-bit weight-only kernel; a packed 4-bit linear dequantises
on every forward pass, so timing it would measure the unpacking kernel rather than the
compression. We therefore report $b{=}4$ for quality and size but exclude it from latency
tables, and answer the sparsity-to-speedup question from the pruning-only arm, whose weights
remain FP32 and benchmark natively.

\paragraph{Never infer speedup from sparsity.} We report three quantities separately:
theoretical non-zero parameter reduction, measured serialised checkpoint size, and measured
CPU latency. An unstructured mask stores zeros, and a dense GEMM kernel multiplies by them
rather than skipping them; a null latency result under unstructured sparsity is a finding
about the deployment path, not a failure of the experiment.

\subsection{Post-hoc recovery ablation}
\label{subsec:recovery}

The core method performs no fine-tuning. Separately, and after the confirmatory result was
known, we examine whether a short full-model recovery phase applied \emph{equally} to both
arms closes or widens the gap. Masks are frozen, fake quantisation remains active, and both
arms receive an identical step count, token budget and learning rate on a recovery slice
verified disjoint from the calibration set. This ablation is labelled exploratory throughout:
it is post-hoc, runs on the validation split, and cannot alter or reinterpret the
confirmatory result. We report it because a recovery phase changes which pipeline is
preferable, and because the magnitude of the recovery effect relative to the pipeline effect
is itself informative for practitioners.

% =====================================================================
%  NOTE FOR CO-AUTHORS -- remove before submission.
%
%  Two repository inconsistencies were encountered while drafting this section and should be
%  resolved before the manuscript is circulated:
%
%  1. docs/method_definition.md ends with a "Current Status" table stating that module
%     selection, the layerwise driver, and all five arms are placeholders. That table is
%     stale: layerwise.py implements the driver and all five arms, and the confirmatory
%     grid (F-37) has been executed. Per CLAUDE.md ("if code and docs disagree, that is a
%     bug in one of them"), the table needs updating rather than silent divergence.
%
%  2. docs/method_definition.md still specifies 50% / 70% sparsity budgets in places, while
%     the frozen confirmatory configuration uses 30% sparsity at both budgets, varying
%     precision (W8 / W4) instead. The frozen values in main_scale_sweep.yaml are the ones
%     used above; method_definition.md should be reconciled to match.
% =====================================================================
